% GENERATED FILE. Edit canonical lessons, then run npm run generate:books.
% canonical-source-hash: fnv1a64:c364f55c1ca21487
% canonical-lessons: ZH-C55-xiang, ZH-W55-xiang, ZH-C55-ai, ZH-C55-dianying, ZH-C55-changge, ZH-C55-kafei

\chapter{What You Love}
\label{ch:zh-what-you-love}
\hlchaptermodality{55}

\begin{chapteropening}
\textbf{By the end of this chapter:} \emph{I can say what I would like to do and what I love.}

Complete the last lesson of chapter 55: I can say what I would like to do and what I love.
\end{chapteropening}

\section[\texorpdfstring{\zh{想} (xiǎng)}{想 (xiǎng)}]{\zh{想} (xiǎng) — to want to; to think; to miss}

\label{lesson:ZH-C55-xiang}



\begin{quote}
Before the new one: say the Mandarin for to ask, then the Mandarin for lesson.
\end{quote}

\subsection*{You'll want to know: \zh{想}}
\textbf{\zh{想}} — \emph{xiǎng} — \textquotedblleft{}to want to; to think; to miss\textquotedblright{}.

\begin{quote}
\textbf{\zh{我想说中国话}} \emph{wǒ xiǎng shuō Zhōngguó huà} — I'd like to speak Chinese.
\end{quote}

The new character, in pieces: \zh{想} — a tree piece and an eye piece side by side, over a heart piece. The next lesson puts it in your hand.

\emph{xiǎng} before another verb says what you would like to do: \emph{wǒ xiǎng shuō}.

\begin{cousinweb}
\textbf{\zh{想}} (\emph{xiǎng}) is inherited from Middle Chinese \emph{sjangX} (Guangyun, in Baxter's transcription), glossed 'to think'.
\end{cousinweb}

\begin{grammarlens}[title={What you've built}]
One more word for everyday answers.
\end{grammarlens}

\subsection*{Your turn}
\begin{itemize}
\raggedright
  \item \emph{Say it:} \emph{xiǎng}
  \item \emph{Say it:} \emph{xiǎng}, once more
  \item \emph{Say it:} \emph{wǒ xiǎng shuō Zhōngguó huà} — I'd like to speak Chinese.
  \item \emph{Recall:} say the Mandarin for to ask, then the Mandarin for lesson, then say \emph{xiǎng} again
\end{itemize}

\begin{tcolorbox}[breakable,colback=teal!4,colframe=teal!35!black,title={Before you move on}]
What does \zh{想} mean? (\textquotedblleft{}To want to; to think; to miss\textquotedblright{}.) Say it once more.
\end{tcolorbox}
\section[\texorpdfstring{\zh{想} (xiǎng)}{想 (xiǎng)}]{\zh{想} — a tree and an eye, over a heart}

\label{lesson:ZH-W55-xiang}



\begin{quote}
The word first, then the character inside it.

\begin{itemize}
\raggedright
  \item \emph{Recall:} say \emph{xiǎng} — \textbf{R1}, one lesson back, before it decays
\end{itemize}

\begin{quote}
\zh{看} \zh{说}
\end{quote}

\emph{Read it:} each one aloud, with nothing said first — \emph{kàn}, \emph{shuō}
\end{quote}

\subsection*{Script — a tree and an eye, over a heart}
\hlblockfigure{\detokenize{figures/ZH-W55-xiang-filmstrip.pdf}}{How \zh{想} is written, stroke by stroke}

\begin{quote}
\zh{想}
\end{quote}

\textbf{Three pieces: the tree, then the eye, then the heart below.} Count three, not thirteen.

\textbf{The tree, four strokes}, narrow on the top left: a horizontal, a vertical through it, a left-falling stroke, and a dot.

\textbf{The eye, five strokes}, on the top right: a vertical, a horizontal that turns down, and three short horizontals inside, the last one closing it. It is the box at the bottom of \zh{看}.

\textbf{The heart, four strokes}, wide underneath: a dot, a long curve that sweeps right and hooks up, then two more dots.

4 + 5 + 4 = thirteen.

\subsection*{Your turn}
\begin{itemize}
\raggedright
  \item \emph{Write it:} \zh{想}, all thirteen, with the model in front of you
  \item \emph{Cover:} write it again from memory, and count the pieces, not the strokes
\end{itemize}

\begin{tcolorbox}[breakable,colback=teal!4,colframe=teal!35!black,title={Before you move on}]
How many pieces in \zh{想}? (\textbf{Three} — the tree, the eye, the heart.) Which character you know has the same eye in it? (\textbf{\zh{看}}.) What does the heart begin with? (\textbf{A dot}.)

Source: \href{https://raw.githubusercontent.com/chanind/hanzi-writer-data/68d10a4b21150cae5e1ebbd223eed289cf32d90c/data/\%E6\%83\%B3.json}{Hanzi Writer Data \zh{想}, pinned stroke paths and medians}.
\end{tcolorbox}
\section[ài]{ài — to love}

\label{lesson:ZH-C55-ai}



\begin{quote}
Before the new one: say the Mandarin for lesson, then the Mandarin for to want to.
\end{quote}

\subsection*{You'll want to know: ài}
\emph{ài} — \textquotedblleft{}to love\textquotedblright{}.

\begin{quote}
\emph{wǒ ài nǐ} — I love you.
\end{quote}

Its characters come later; for now it is a word for the ear.

\begin{cousinweb}
\emph{ài} is inherited from Middle Chinese \emph{'ojH} (Guangyun, in Baxter's transcription), glossed 'to cherish'.
\end{cousinweb}

\begin{grammarlens}[title={What you've built}]
One more word for everyday answers.
\end{grammarlens}

\subsection*{Your turn}
\begin{itemize}
\raggedright
  \item \emph{Say it:} \emph{ài}
  \item \emph{Say it:} \emph{ài}, once more
  \item \emph{Say it:} \emph{wǒ ài nǐ} — I love you.
  \item \emph{Recall:} say the Mandarin for lesson, then the Mandarin for to want to, then say \emph{ài} again
\end{itemize}

\begin{tcolorbox}[breakable,colback=teal!4,colframe=teal!35!black,title={Before you move on}]
What does ài mean? (\textquotedblleft{}To love\textquotedblright{}.) Say it once more.
\end{tcolorbox}
\section[diànyǐng]{diànyǐng — film; movie}

\label{lesson:ZH-C55-dianying}



\begin{quote}
Before the new one: say the Mandarin for to want to, then the Mandarin for to love.
\end{quote}

\subsection*{You'll want to know: diànyǐng}
\emph{diànyǐng} — \textquotedblleft{}film; movie\textquotedblright{}.

\begin{quote}
\emph{wǒ xǐhuan kàn diànyǐng} — I like going to the cinema.
\end{quote}

Its characters come later; for now it is a word for the ear.

\begin{cousinweb}
\emph{diànyǐng} is built from diàn 'lightning', now 'electricity', and yǐng 'shadow', both inherited from Middle Chinese, \emph{denH} and \emph{'jaengX} (Guangyun, in Baxter's transcription): electric shadows.
\end{cousinweb}

\begin{grammarlens}[title={What you've built}]
One more word for everyday answers.
\end{grammarlens}

\subsection*{Your turn}
\begin{itemize}
\raggedright
  \item \emph{Say it:} \emph{diànyǐng}
  \item \emph{Say it:} \emph{diànyǐng}, once more
  \item \emph{Say it:} \emph{wǒ xǐhuan kàn diànyǐng} — I like going to the cinema.
  \item \emph{Recall:} say the Mandarin for to want to, then the Mandarin for to love, then say \emph{diànyǐng} again
\end{itemize}

\begin{tcolorbox}[breakable,colback=teal!4,colframe=teal!35!black,title={Before you move on}]
What does diànyǐng mean? (\textquotedblleft{}Film; movie\textquotedblright{}.) Say it once more.
\end{tcolorbox}
\section[chànggē]{chànggē — to sing}

\label{lesson:ZH-C55-changge}



\begin{quote}
Before the new one: say the Mandarin for to love, then the Mandarin for film.
\end{quote}

\subsection*{You'll want to know: chànggē}
\emph{chànggē} — \textquotedblleft{}to sing\textquotedblright{}.

\begin{quote}
\emph{tā xǐhuan chànggē} — She loves singing.
\end{quote}

Its characters come later; for now it is a word for the ear.

\begin{cousinweb}
\emph{chànggē} is built from chàng 'to strike up a song' and gē 'song', both inherited from Middle Chinese, \emph{tsyhangH} and \emph{ka} (Guangyun, in Baxter's transcription).
\end{cousinweb}

\begin{grammarlens}[title={What you've built}]
One more word for everyday answers.
\end{grammarlens}

\subsection*{Your turn}
\begin{itemize}
\raggedright
  \item \emph{Say it:} \emph{chànggē}
  \item \emph{Say it:} \emph{chànggē}, once more
  \item \emph{Say it:} \emph{tā xǐhuan chànggē} — She loves singing.
  \item \emph{Recall:} say the Mandarin for to love, then the Mandarin for film, then say \emph{chànggē} again
\end{itemize}

\begin{tcolorbox}[breakable,colback=teal!4,colframe=teal!35!black,title={Before you move on}]
What does chànggē mean? (\textquotedblleft{}To sing\textquotedblright{}.) Say it once more.
\end{tcolorbox}
\section[kāfēi]{kāfēi — coffee}

\label{lesson:ZH-C55-kafei}



\begin{quote}
Before the new one: say the Mandarin for film, then the Mandarin for to sing.
\end{quote}

\subsection*{You'll want to know: kāfēi}
\emph{kāfēi} — \textquotedblleft{}coffee\textquotedblright{}.

\begin{quote}
\emph{nǐ hē kāfēi ma} — Do you drink coffee?
\end{quote}

Its characters come later; for now it is a word for the ear.

\begin{cousinweb}
\emph{kāfēi} is borrowed from a European word for coffee; it is usually traced to English \emph{coffee}, though the path is not documented in what we can open.
\end{cousinweb}

\begin{grammarlens}[title={What you've built}]
One more word for everyday answers.
\end{grammarlens}

\subsection*{Your turn}
\begin{itemize}
\raggedright
  \item \emph{Say it:} \emph{kāfēi}
  \item \emph{Say it:} \emph{kāfēi}, once more
  \item \emph{Say it:} \emph{nǐ hē kāfēi ma} — Do you drink coffee?
  \item \emph{Recall:} say the Mandarin for film, then the Mandarin for to sing, then say \emph{kāfēi} again
\end{itemize}

\begin{tcolorbox}[breakable,colback=teal!4,colframe=teal!35!black,title={Before you move on}]
What does kāfēi mean? (\textquotedblleft{}Coffee\textquotedblright{}.) Say it once more.
\end{tcolorbox}
